pdfoutput=1
% Options for packages loaded elsewhere
\PassOptionsToPackage{unicode}{hyperref}
\PassOptionsToPackage{hyphens}{url}
\pdfoutput=1
\documentclass[
  12pt,
]{article}
\usepackage{xcolor}
\usepackage{amsmath,amssymb}
\setcounter{secnumdepth}{-\maxdimen} % remove section numbering
\usepackage{iftex}
\ifPDFTeX
  \usepackage[T1]{fontenc}
  \usepackage{textcomp} % provide euro and other symbols
\else % if luatex or xetex
  \usepackage{unicode-math} % this also loads fontspec
  \defaultfontfeatures{Scale=MatchLowercase}
  \defaultfontfeatures[\rmfamily]{Ligatures=TeX,Scale=1}
\fi
\usepackage{lmodern}
\ifPDFTeX\else
  % xetex/luatex font selection
\fi
% Use upquote if available, for straight quotes in verbatim environments
\IfFileExists{upquote.sty}{\usepackage{upquote}}{}
\IfFileExists{microtype.sty}{% use microtype if available
  \usepackage[]{microtype}
  \UseMicrotypeSet[protrusion]{basicmath} % disable protrusion for tt fonts
}{}
\makeatletter
\@ifundefined{KOMAClassName}{% if non-KOMA class
  \IfFileExists{parskip.sty}{%
    \usepackage{parskip}
  }{% else
    \setlength{\parindent}{0pt}
    \setlength{\parskip}{6pt plus 2pt minus 1pt}}
}{% if KOMA class
  \KOMAoptions{parskip=half}}
\makeatother
\usepackage{longtable,booktabs,array}
\usepackage{calc} % for calculating minipage widths
% Correct order of tables after \paragraph or \subparagraph
\usepackage{etoolbox}
\makeatletter
\patchcmd\longtable{\par}{\if@noskipsec\mbox{}\fi\par}{}{}
\makeatother
% Allow footnotes in longtable head/foot
\IfFileExists{footnotehyper.sty}{\usepackage{footnotehyper}}{\usepackage{footnote}}
\makesavenoteenv{longtable}
\setlength{\emergencystretch}{3em} % prevent overfull lines
\providecommand{\tightlist}{%
  \setlength{\itemsep}{0pt}\setlength{\parskip}{0pt}}
\usepackage{bookmark}
\IfFileExists{xurl.sty}{\usepackage{xurl}}{} % add URL line breaks if available
\urlstyle{same}
\hypersetup{
  hidelinks,
  pdfcreator={LaTeX via pandoc}}

\author{SCX}
\date{2026}

\begin{document}

\section{Review: Minimax Lower Bound v2 (Hellinger Distance
Proof)}\label{review-minimax-lower-bound-v2-hellinger-distance-proof}

\begin{quote}
\textbf{File reviewed}:
\texttt{G:\textbackslash{}Xiaogan\_Supercomputing\_data\textbackslash{}SCX\textbackslash{}theory\textbackslash{}explorations\textbackslash{}minimax\_lower\_bound\_v2.md}

\textbf{Review date}: 2026-06-27

\textbf{Scope}: Full mathematical verification of all inequalities,
formulas, and derivations. Numerical cross-checking of all claimed
relationships.
\end{quote}

\begin{center}\rule{0.5\linewidth}{0.5pt}\end{center}

\subsection{1. Overall Verdict}\label{overall-verdict}

The proof correctly fixes the three fatal issues from v1 (Slud's
inequality, chi-squared direction, F1 algebra), and the core Hellinger
tensorization approach is sound. \textbf{However, two new fatal errors
have been introduced in the F1 conversion section} (Section 7), and a
significant formula error appears in the K \textgreater{} 2 section
(Section 8.3).

\textbf{Bottom line}: The testing error lower bound (Part a) is correct.
The F1 lower bound (Part b) has two algebraic errors that invalidate its
stated constant. The rate optimality claim (Part c) remains correct
because the exponential rate \texttt{exp(-2MDelta\^{}2)} is preserved
under the corrected analysis.

\begin{center}\rule{0.5\linewidth}{0.5pt}\end{center}

\subsection{2. Issues Found}\label{issues-found}

\subsubsection{Issue 1 {[}FATAL{]}: Section 7.3 -- Backwards inequality
in F1
simplification}\label{issue-1-fatal-section-7.3-backwards-inequality-in-f1-simplification}

\textbf{Location}: Lines 423-427

\textbf{Claim}:

\begin{verbatim}
(1-eta)*(rho^M/4)/(2eta+(1-eta)*(rho^M/4)) >= (1-eta)*rho^M/(8eta)
\end{verbatim}

\textbf{What is wrong}: The inequality is backwards. Since
\texttt{2eta\ +\ (1-eta)*(rho\^{}M/4)\ \textgreater{}\ 2eta} for any
positive \texttt{rho\^{}M}, the left-hand side has a larger denominator
with the same numerator, making it strictly smaller. The correct
inequality is \texttt{\textless{}=}, not \texttt{\textgreater{}=}.

\textbf{Verification} (eta = 0.3, rho\^{}M = 0.1):

\begin{verbatim}
LHS = 0.7*0.025/(0.6+0.7*0.025) = 0.02834
RHS = 0.7*0.1/(8*0.3) = 0.02917
LHS < RHS, so LHS >= RHS is FALSE.
\end{verbatim}

\textbf{Consequence}: The simplified bound
\texttt{1-F1\ \textgreater{}=\ rho\^{}M/(16eta)} is not validly derived.
The correct inequality direction gives an upper bound
\texttt{LHS\ \textless{}=\ RHS}, which is useless for lower-bounding
1-F1.

\textbf{How to fix}: The exact bound must be kept in its unsimplified
form:

\begin{verbatim}
1-F1 >= (1-eta)(rho^M/4)/(2eta+(1-eta)(rho^M/4))
\end{verbatim}

Or a correct simplification with a multiplicative correction:

\begin{verbatim}
1-F1 >= (1-eta)*rho^M/(8eta) * (1 - (1-eta)*rho^M/(4eta))
\end{verbatim}

(using \texttt{1/(1+x)\ \textgreater{}=\ 1-x} for
\texttt{x\ \textgreater{}=\ 0})

\begin{center}\rule{0.5\linewidth}{0.5pt}\end{center}

\subsubsection{Issue 2 {[}FATAL{]}: Section 7.5 -- Wrong case selected
as tighter
bound}\label{issue-2-fatal-section-7.5-wrong-case-selected-as-tighter-bound}

\textbf{Location}: Lines 439-467

\textbf{Claim}:

\begin{verbatim}
eps/2 >= (1-eta)*eps/(2eta+(1-eta)*eps) for eta <= 1/2
\end{verbatim}

And consequently:

\begin{verbatim}
min(eps/(2-eps), (1-eta)*eps/(2eta+(1-eta)*eps)) = (1-eta)*eps/(2eta+(1-eta)*eps)
\end{verbatim}

\textbf{What is wrong}: The inequality direction is reversed for
\texttt{eta\ \textless{}\ 1/2} and small \texttt{eps} (the regime of
interest). The proof attempts to verify this by cross-multiplication,
claiming:

\texttt{2eta\ +\ (1-eta)*eps\ \textgreater{}=\ 2*(1-eta)} follows from
\texttt{eps\ \textgreater{}=\ 0,\ eta\ \textgreater{}=\ 0}.

But the correct cross-multiplication is:

\begin{verbatim}
eps/2 >= (1-eta)*eps/(2eta+(1-eta)*eps)
=> 2eta + (1-eta)*eps >= 2*(1-eta)
=> (1-eta)*eps >= 2 - 4eta
\end{verbatim}

For \texttt{eta\ =\ 0.3}: requires
\texttt{eps\ \textgreater{}=\ 0.8/0.7\ ≈\ 1.14}. Impossible since
\texttt{eps\ \textless{}=\ 1} (and typically
\texttt{eps\ \textless{}\textless{}\ 1}). For \texttt{eta\ =\ 0.4}:
requires \texttt{eps\ \textgreater{}=\ 0.4/0.6\ ≈\ 0.667}. Possible but
atypical -- rho\^{}M would need to be \textgreater= 2.67, impossible
since rho \textless= 1. For \texttt{eta\ =\ 0.49}: requires
\texttt{eps\ \textgreater{}=\ 0.04/0.51\ ≈\ 0.078}. Possible for some
parameters. For \texttt{eta\ =\ 0.5}: requires
\texttt{eps\ \textgreater{}=\ 0}. True.

So the claim only holds when
\texttt{(1-eta)*eps\ \textgreater{}=\ 2-4eta}, i.e., when
\texttt{eta\ \textgreater{}=\ 1/(2-eps/(1-eps))} approximately. For most
of the parameter space (especially small \texttt{eps} and
small-to-moderate \texttt{eta}), the claim is FALSE.

\textbf{Verification} (eta = 0.3, eps = 0.1):

\begin{verbatim}
eps/2 = 0.05
(1-eta)*eps/(2eta+(1-eta)*eps) = 0.7*0.1/0.67 ≈ 0.104
So eps/2 < case2_bound, not >=.
\end{verbatim}

The correct minimum is:
\texttt{min(eps/(2-eps),\ (1-eta)*eps/(2eta+(1-eta)*eps))\ =\ eps/(2-eps)}
(Case 1)

for all \texttt{eta\ \textless{}\ 1/2} and small \texttt{eps}.

\textbf{Consequence}: The claimed universal bound
\texttt{1-F1\ \textgreater{}=\ (1-eta)*eps/(2eta+(1-eta)*eps)} is
invalid. The correct universal bound uses Case 1:

\begin{verbatim}
1-F1 >= eps/(2-eps) = rho^M/(8-rho^M)
\end{verbatim}

which is approximately \texttt{rho\^{}M/8} (not
\texttt{rho\^{}M/(16eta)}).

\textbf{Comparison of bounds}:

\begin{verbatim}
eta=0.1: correct=rho^M/8, claimed=rho^M/1.6  (claimed 5x too large)
eta=0.3: correct=rho^M/8, claimed=rho^M/4.8  (claimed 1.67x too large)
eta=0.5: correct=rho^M/8, claimed=rho^M/8    (match at eta=0.5)
\end{verbatim}

\textbf{How to fix}: Use the Case 1 bound:

\begin{verbatim}
1-F1 >= eps/(2-eps) = (rho^M/4)/(2-rho^M/4) = rho^M/(8-rho^M)
\end{verbatim}

For the simplified form (small rho\^{}M):

\begin{verbatim}
1-F1 >= rho^M/8
\end{verbatim}

\begin{center}\rule{0.5\linewidth}{0.5pt}\end{center}

\subsubsection{Issue 3 {[}MAJOR{]}: Section 8.3 -- Incorrect Hellinger
affinity formula for K \textgreater{}
2}\label{issue-3-major-section-8.3-incorrect-hellinger-affinity-formula-for-k-2}

\textbf{Location}: Line 511

\textbf{Claim}:

\begin{verbatim}
rho_K = 2*sqrt(mu_s * (1 - mu_s/(K-1)))
\end{verbatim}

\textbf{What is wrong}: This formula is only correct for K=2. The
general Hellinger affinity for \texttt{Bernoulli(p)} vs
\texttt{Bernoulli(q)} is:

\begin{verbatim}
rho = sqrt(p*q) + sqrt((1-p)*(1-q))
\end{verbatim}

The proof's formula \texttt{2*sqrt(p*q)} only matches the general
formula when \texttt{q\ =\ 1-p}, i.e., when \texttt{K=2} (since
\texttt{1\ -\ mu\_s/(K-1)\ =\ 1\ -\ mu\_s\ =\ 1\ -\ p} only when
\texttt{K=2}).

For K \textgreater{} 2, the correct formula is:

\begin{verbatim}
rho_K = sqrt(mu_s * (1 - mu_s/(K-1))) + sqrt((1-mu_s) * (mu_s/(K-1)))
\end{verbatim}

\textbf{Verification} (mu\_s = 0.3, K = 3):

\begin{verbatim}
correct: sqrt(0.3*0.85) + sqrt(0.7*0.15) = 0.505 + 0.324 = 0.829
wrong:   2*sqrt(0.3*0.85) = 2*0.505 = 1.010
\end{verbatim}

The wrong formula gives \texttt{rho\ \textgreater{}\ 1} for many
parameter values (e.g., mu=0.3, K=3: rho=1.010; mu=0.2, K=5: rho=0.872).
While \texttt{rho\ \textless{}=\ 1} should always hold for Hellinger
affinity. The correct formula always satisfies
\texttt{rho\ \textless{}=\ 1}.

\textbf{Consequence}: The numerical values and asymptotic expressions
for K \textgreater{} 2 are wrong. However, the \textbf{conclusion} that
K=2 is the hardest case is still correct (verified with the correct
formula -- rho\_K \textless{} rho\_2 for K \textgreater{} 2, giving a
weaker bound, so K=2 gives the strongest valid lower bound for all K
\textgreater= 2).

\textbf{How to fix}: Replace the formula with the correct expression:

\begin{verbatim}
rho_K = sqrt(mu_s * (1 - mu_s/(K-1))) + sqrt((1-mu_s) * (mu_s/(K-1)))
\end{verbatim}

And note that
\texttt{rho\_K\ \textless{}\ rho\_2\ =\ 2*sqrt(mu\_s*(1-mu\_s))} for all
K \textgreater{} 2, which can be verified by squaring both sides.

\begin{center}\rule{0.5\linewidth}{0.5pt}\end{center}

\subsubsection{Issue 4 {[}MAJOR{]}: Section 11 -- Chernoff exactness
claim fails for K \textgreater{}
2}\label{issue-4-major-section-11-chernoff-exactness-claim-fails-for-k-2}

\textbf{Location}: Lines 676-681

\textbf{Claim}:

\begin{verbatim}
-log(rho) = KL(1/2 || mu_s) = Chernoff exponent
\end{verbatim}

\textbf{What is wrong}: This holds for K=2 because the likelihood ratio
\texttt{dP1/dP0} is symmetric on the log scale (values are reciprocals).
For K \textgreater{} 2, the optimal Chernoff exponent generally uses a
tilting parameter \texttt{t\ !=\ 1/2}, so
\texttt{-log(rho\_K)\ \textgreater{}\ KL(1/2\ \textbar{}\textbar{}\ mu\_s)}.
The Hellinger exponent \texttt{-log(rho\_K)} is still a valid lower
bound on the Chernoff exponent, but it is \textbf{not exact}.

\textbf{Verification} (mu=0.3, K=3):

\begin{verbatim}
-log(rho_K) = -log(0.829) = 0.188
KL(1/2||0.3) = 0.087
-0.188 != 0.087.
\end{verbatim}

\textbf{Consequence}: The claim that the Hellinger approach gives
``exact Chernoff exponent'' is true only for K=2 (or more generally,
when the two distributions are symmetric on the log-odds scale). For K
\textgreater{} 2, the Hellinger exponent is a valid but non-exact lower
bound.

\textbf{How to fix}: Clarify that the Chernoff exactness holds for K=2
and provide the general expression for K \textgreater{} 2, or simply
note that the Hellinger bound is always valid and the exponent is
bounded below by the K=2 exponent.

\begin{center}\rule{0.5\linewidth}{0.5pt}\end{center}

\subsubsection{Issue 5 {[}MAJOR{]}: Section 9 -- C\_bal \textgreater{} 1
extension is
incomplete}\label{issue-5-major-section-9-c_bal-1-extension-is-incomplete}

\textbf{Location}: Lines 536-619

\textbf{What is wrong}: The section is a sketch, not a rigorous proof.
There are several gaps:

\begin{enumerate}
\def\labelenumi{\arabic{enumi}.}
\item
  The convexity bound via \texttt{sqrt} concavity gives
  \texttt{H\^{}2(P0,\ P1)\ \textless{}=\ avg\ H\^{}2(P0,\ Q\_ell)},
  which is correct. But the translation to a lower bound on the testing
  error uses \texttt{TV\ \textless{}=\ H}, which yields
  \texttt{R\ \textgreater{}=\ (1-TV)/2\ \textgreater{}=\ (1-sqrt(avg\ H\^{}2))/2}.
  This is a valid bound but may be very weak.
\item
  The analysis switches between using \texttt{max\ H} and
  \texttt{min\ H} inconsistently. Line 605 says ``the worst (largest)
  Hellinger distance among the mixture components gives the tightest
  upper bound on TV'', but line 610 then says ``the component with the
  smallest\ldots{} gives the smallest Hellinger distance, and hence the
  weakest lower bound.'' These are different criteria and the
  presentation is confusing.
\item
  The claim ``the exponent 2MDelta\^{}2 is preserved'' (line 618) is
  asserted without proof. While it's plausibly true (each component's
  gap is at least \texttt{Delta} times some C\_bal-dependent factor), no
  rigorous derivation is given.
\end{enumerate}

\textbf{Consequence}: The Theorem 4 statement claims to hold under
(A1)-(A6) without requiring C\_bal = 1, but the proof only fully handles
C\_bal = 1. The C\_bal \textgreater{} 1 case is not rigorously proven.

\textbf{How to fix}: Either (a) add C\_bal = 1 as an explicit condition
in the theorem statement, or (b) provide a complete derivation for
C\_bal \textgreater{} 1. Option (b) would involve bounding the average
Hellinger affinity:

\begin{verbatim}
(1/L) sum_ell rho_ell^M >= min_ell rho_ell^M
\end{verbatim}

where
\texttt{rho\_ell\ =\ sqrt(mu\_s\ *\ p\_ell)\ +\ sqrt((1-mu\_s)\ *\ (1-p\_ell))}
and \texttt{p\_ell\ =\ 1\ -\ C\_ell*mu\_s/(K-1)} for some
\texttt{C\_ell\ \textgreater{}=\ 1}. The smallest \texttt{rho\_ell}
(closest to 1) gives the weakest bound, and this is attained at
\texttt{C\_ell\ =\ C\_bal} (or similar).

\begin{center}\rule{0.5\linewidth}{0.5pt}\end{center}

\subsubsection{Issue 6 {[}MINOR{]}: Section 6.4 -- Wrong Taylor
coefficient for
Delta\^{}4}\label{issue-6-minor-section-6.4-wrong-taylor-coefficient-for-delta4}

\textbf{Location}: Lines 340-344

\textbf{Claim}:

\begin{verbatim}
log(1-4Delta^2) = -4Delta^2 - 8Delta^4/3 - O(Delta^6)
\end{verbatim}

\textbf{What is wrong}: The Taylor expansion of
\texttt{log(1-x)\ =\ -x\ -\ x\^{}2/2\ -\ x\^{}3/3\ -\ ...} with
\texttt{x\ =\ 4Delta\^{}2} gives:

\begin{verbatim}
log(1-4Delta^2) = -4Delta^2 - (4Delta^2)^2/2 - (4Delta^2)^3/3 - ...
               = -4Delta^2 - 16Delta^4/2 - 64Delta^6/3 - ...
               = -4Delta^2 - 8Delta^4 - O(Delta^6)
\end{verbatim}

The coefficient of \texttt{Delta\^{}4} is \texttt{-8}, not
\texttt{-8/3}.

Then
\texttt{(M/2)*log(1-4Delta\^{}2)\ =\ -2M*Delta\^{}2\ -\ 4M*Delta\^{}4\ -\ O(M*Delta\^{}6)},
not \texttt{-2M*Delta\^{}2\ -\ (4/3)M*Delta\^{}4\ -\ O(M*Delta\^{}6)}.

\textbf{Consequence}: The specific coefficient in the
\texttt{O(M*Delta\^{}4)} term is wrong, but this term is absorbed into
the \texttt{O(M*Delta\^{}4)} notation anyway. No substantive claim is
affected.

\textbf{How to fix}: Correct the coefficient.

\begin{center}\rule{0.5\linewidth}{0.5pt}\end{center}

\subsection{3. Algebraic Verification of F1
Conversion}\label{algebraic-verification-of-f1-conversion}

This is the most critical section. Below is a complete, corrected
derivation of the F1 lower bound.

\subsubsection{Setup}\label{setup}

\begin{itemize}
\tightlist
\item
  Testing error: \texttt{R(psi)\ =\ max(alpha,\ beta)} where
  \texttt{alpha\ =\ P0(psi=1)}, \texttt{beta\ =\ P1(psi=0)}
\item
  Le Cam bound:
  \texttt{R(psi)\ \textgreater{}=\ (1-TV)/2\ \textgreater{}=\ rho\^{}M/4\ =:\ epsilon}
\item
  So \texttt{max(alpha,\ beta)\ \textgreater{}=\ epsilon}
\end{itemize}

\subsubsection{F1 in terms of alpha,
beta}\label{f1-in-terms-of-alpha-beta}

\begin{verbatim}
F1 = 2*TP/(2*TP + FP + FN)
   = 2*eta*(1-beta) / (2*eta*(1-beta) + (1-eta)*alpha + eta*beta)
\end{verbatim}

\subsubsection{Case analysis}\label{case-analysis}

\textbf{Case 1}: \texttt{beta\ \textgreater{}=\ epsilon} (false negative
rate is high)

Using \texttt{FP\ \textgreater{}=\ 0}:

\begin{verbatim}
F1 <= 2*eta*(1-beta) / (2*eta*(1-beta) + eta*beta)
    = 2*(1-beta)/(2-beta)
\end{verbatim}

Since \texttt{beta\ \textgreater{}=\ epsilon} and the function is
decreasing in \texttt{beta}:

\begin{verbatim}
F1 <= 2*(1-epsilon)/(2-epsilon)
1-F1 >= epsilon/(2-epsilon)
\end{verbatim}

\textbf{Case 2}: \texttt{alpha\ \textgreater{}=\ epsilon} (false
positive rate is high)

Using \texttt{FN\ \textgreater{}=\ 0}:

\begin{verbatim}
F1 <= 2*eta*(1-beta) / (2*eta*(1-beta) + (1-eta)*alpha)
\end{verbatim}

Since \texttt{beta\ \textgreater{}=\ 0} (maximizes F1 at
\texttt{beta\ =\ 0}) and \texttt{alpha\ \textgreater{}=\ epsilon}:

\begin{verbatim}
F1 <= 2*eta / (2*eta + (1-eta)*epsilon)
1-F1 >= (1-eta)*epsilon / (2*eta + (1-eta)*epsilon)
\end{verbatim}

\subsubsection{Universal bound}\label{universal-bound}

Since \texttt{max(alpha,\ beta)\ \textgreater{}=\ epsilon}, either Case
1 or Case 2 must hold. Therefore:

\begin{verbatim}
1-F1 >= min(epsilon/(2-epsilon), (1-eta)*epsilon/(2*eta+(1-eta)*epsilon))
\end{verbatim}

\subsubsection{Which case is tighter?}\label{which-case-is-tighter}

For \texttt{eta\ \textless{}\ 1/2} and small \texttt{epsilon} (the
typical regime):

\begin{verbatim}
epsilon/(2-epsilon) ≈ epsilon/2
(1-eta)*epsilon/(2*eta+(1-eta)*epsilon) ≈ (1-eta)*epsilon/(2*eta)
\end{verbatim}

Ratio:
\texttt{(epsilon/2)\ /\ ((1-eta)*epsilon/(2*eta))\ =\ eta/(1-eta)\ \textless{}\ 1}
for \texttt{eta\ \textless{}\ 1/2}.

So Case 1 gives the smaller bound (tighter). \textbf{This contradicts
the proof's claim.}

\subsubsection{Correct F1 lower bound}\label{correct-f1-lower-bound}

\begin{verbatim}
1-F1 >= epsilon/(2-epsilon) = (rho^M/4)/(2 - rho^M/4) = rho^M/(8 - rho^M)
\end{verbatim}

For small \texttt{rho\^{}M}:

\begin{verbatim}
1-F1 >= rho^M/8   (conservative simplification)
\end{verbatim}

Compare with the claimed bound \texttt{rho\^{}M/(16*eta)}:

\begin{longtable}[]{@{}llll@{}}
\toprule\noalign{}
eta & Correct bound & Claimed bound & Ratio \\
\midrule\noalign{}
\endhead
\bottomrule\noalign{}
\endlastfoot
0.1 & \texttt{rho\^{}M/8} & \texttt{rho\^{}M/1.6} & 5x too large \\
0.3 & \texttt{rho\^{}M/8} & \texttt{rho\^{}M/4.8} & 1.67x too large \\
0.5 & \texttt{rho\^{}M/8} & \texttt{rho\^{}M/8} & Match \\
\end{longtable}

The claimed bound is \textbf{stronger than justified} for all
\texttt{eta\ \textless{}\ 1/2}, with the discrepancy growing as
\texttt{eta} decreases.

\subsubsection{Rate optimality
preserved}\label{rate-optimality-preserved}

Both the correct bound \texttt{rho\^{}M/8} and the claimed bound
\texttt{rho\^{}M/(16*eta)} have leading term \texttt{rho\^{}M} =
\texttt{(2*sqrt(mu\_s(1-mu\_s)))\^{}M} =
\texttt{exp(-2M*Delta\^{}2\ +\ O(M*Delta\^{}4))}. The rate is preserved,
confirming that Part (c) is unaffected.

\begin{center}\rule{0.5\linewidth}{0.5pt}\end{center}

\subsection{4. Summary of Correct vs.~Claimed
Results}\label{summary-of-correct-vs.-claimed-results}

\begin{longtable}[]{@{}
  >{\raggedright\arraybackslash}p{(\linewidth - 6\tabcolsep) * \real{0.2000}}
  >{\raggedright\arraybackslash}p{(\linewidth - 6\tabcolsep) * \real{0.3000}}
  >{\raggedright\arraybackslash}p{(\linewidth - 6\tabcolsep) * \real{0.3000}}
  >{\raggedright\arraybackslash}p{(\linewidth - 6\tabcolsep) * \real{0.2000}}@{}}
\toprule\noalign{}
\begin{minipage}[b]{\linewidth}\raggedright
Quantity
\end{minipage} & \begin{minipage}[b]{\linewidth}\raggedright
Claimed in v2
\end{minipage} & \begin{minipage}[b]{\linewidth}\raggedright
Correct value
\end{minipage} & \begin{minipage}[b]{\linewidth}\raggedright
Severity
\end{minipage} \\
\midrule\noalign{}
\endhead
\bottomrule\noalign{}
\endlastfoot
Testing bound & \texttt{rho\^{}M/4} & \texttt{rho\^{}M/4} &
\textbf{OK} \\
F1 bound (const) & \texttt{rho\^{}M/(16*eta)} &
\texttt{rho\^{}M/(8-rho\^{}M)} & \textbf{FATAL (const wrong)} \\
F1 bound (rate) & \texttt{exp(-2M*Delta\^{}2)} &
\texttt{exp(-2M*Delta\^{}2)} & \textbf{OK} \\
K\textgreater2 affinity & \texttt{2*sqrt(mu*(1-mu/(K-1)))} &
\texttt{sqrt(mu*(1-mu/(K-1)))\ +\ sqrt((1-mu)*mu/(K-1))} & \textbf{MAJOR
(formula wrong)} \\
Chernoff exactness & Holds for all K & Holds for K=2 only &
\textbf{MAJOR (overclaimed)} \\
C\_bal \textgreater{} 1 & Claimed but not proven & Needs rigorous
treatment & \textbf{MAJOR (gap)} \\
Taylor coeff & \texttt{-8/3*Delta\^{}4} & \texttt{-8*Delta\^{}4} &
\textbf{MINOR} \\
\end{longtable}

\begin{center}\rule{0.5\linewidth}{0.5pt}\end{center}

\subsection{5. Comparison to v1: What is Fixed, What
Remains}\label{comparison-to-v1-what-is-fixed-what-remains}

\subsubsection{Fixed from v1}\label{fixed-from-v1}

\begin{longtable}[]{@{}
  >{\raggedright\arraybackslash}p{(\linewidth - 6\tabcolsep) * \real{0.3043}}
  >{\raggedright\arraybackslash}p{(\linewidth - 6\tabcolsep) * \real{0.1739}}
  >{\raggedright\arraybackslash}p{(\linewidth - 6\tabcolsep) * \real{0.1739}}
  >{\raggedright\arraybackslash}p{(\linewidth - 6\tabcolsep) * \real{0.3478}}@{}}
\toprule\noalign{}
\begin{minipage}[b]{\linewidth}\raggedright
Issue
\end{minipage} & \begin{minipage}[b]{\linewidth}\raggedright
v1
\end{minipage} & \begin{minipage}[b]{\linewidth}\raggedright
v2
\end{minipage} & \begin{minipage}[b]{\linewidth}\raggedright
Status
\end{minipage} \\
\midrule\noalign{}
\endhead
\bottomrule\noalign{}
\endlastfoot
Slud's inequality & Fatal error (false bound) & Hellinger tensorization
& \textbf{Fixed} \\
Chi-squared direction & Used wrong direction for K\textgreater2 &
Product reduction with correct direction & \textbf{Fixed} \\
F1 algebra & Algebraic error in Lemma 8 & Case analysis (new errors
introduced) & \textbf{Partially fixed, new errors} \\
Odd M ceiling & Unaddressed & Works for all M & \textbf{Fixed} \\
\end{longtable}

\subsubsection{Remaining issues (this
review)}\label{remaining-issues-this-review}

\begin{enumerate}
\def\labelenumi{\arabic{enumi}.}
\item
  \textbf{F1 constant is wrong} (Section 7.3, 7.5) -- The simplification
  uses backwards inequality directions. The correct bound is
  \texttt{1-F1\ \textgreater{}=\ rho\^{}M/(8-rho\^{}M)}, which is
  approximately \texttt{rho\^{}M/8}, not \texttt{rho\^{}M/(16*eta)}.
\item
  \textbf{K \textgreater{} 2 affinity formula is wrong} (Section 8.3) --
  The formula \texttt{2*sqrt(mu*(1-mu/(K-1)))} is incorrect and can give
  values \textgreater{} 1. The correct formula is a sum of two square
  roots.
\item
  \textbf{Chernoff exactness overclaimed} (Section 11) -- True for K=2
  but not K\textgreater2.
\item
  \textbf{C\_bal \textgreater{} 1 is a sketch} (Section 9) -- Not a
  complete proof.
\end{enumerate}

\subsubsection{Recommendations}\label{recommendations}

\begin{enumerate}
\def\labelenumi{\arabic{enumi}.}
\tightlist
\item
  Correct the F1 bound to
  \texttt{1-F1\ \textgreater{}=\ rho\^{}M/(8-rho\^{}M)} (or the
  conservative \texttt{rho\^{}M/8}).
\item
  Correct the K\textgreater2 affinity formula.
\item
  Move the Chernoff exactness claim to the K=2 section only.
\item
  Either require C\_bal = 1 explicitly in the theorem statement, or
  provide a complete C\_bal \textgreater{} 1 proof.
\item
  Verify the corrected F1 bound still matches the numerical examples (it
  does: \texttt{1-F1\ \textgreater{}=\ rho\^{}M/8} gives
  \texttt{0.8\^{}10/8\ ≈\ 0.013}, which is \texttt{\textless{}=} the
  direct bound \texttt{0.027} from the full expression -- this is
  conservative but valid).
\end{enumerate}

\begin{center}\rule{0.5\linewidth}{0.5pt}\end{center}

\emph{End of review.}

\end{document}
